\section{神经网络模型}

为系统评估深度学习方法在入侵检测任务上的表现，本章实现了三种典型神经网络模型：多层感知器（MLP）、一维卷积网络（CNN）和长短期记忆网络（LSTM）。深度学习在入侵检测领域的应用已得到广泛研究~\cite{jain2026optimizing,alashjaee2025deep,naz2025innovative}。所有模型均基于 PyTorch 框架构建与训练~\cite{pytorch}，在二分类（正常 vs 异常）和多分类（5 大攻击类别）两个任务上完成训练与评估。

\subsection{MLP多层感知器}

MLP 作为神经网络的基线模型，其架构设计如下：输入层 20 维（与第 2 章特征工程一致），隐藏层依次为 128 和 64，输出层 2 个神经元对应二分类任务，激活函数统一使用 ReLU，并在隐藏层后加入 Dropout(0.3) 以抑制过拟合。训练策略方面，采用 Adam 优化器（学习率 1e-3）、交叉熵损失函数，最大训练轮数 30，并引入 EarlyStopping（patience=5）与 ReduceLROnPlateau 动态学习率调度以提升收敛稳定性。

在 hold-out 验证集上，MLP 二分类基线取得优秀指标：准确率 0.9891、F1 值 0.9883、AUC 0.9992、召回率 0.9850。但应注意，此验证集是从训练集按 8:2 随机切分的子集，与训练集分布完全一致（IID），因此该高指标反映的是模型在已知分布上的拟合能力，并不能直接代表对未见攻击的泛化能力。

为挖掘进一步潜力，对 hidden\_dims 与 activation 进行了 6 组合网格搜索。最优组合为 (256, 128, 64) + ReLU，调优后验证集指标为：acc=0.9902，f1=0.9894，AUC=0.9992。相对基线的边际提升不足 0.2 个百分点，说明验证集指标已接近该数据划分下的上限，调优空间有限。

多分类场景下，准确率 81.0\%、F1\_macro=71.8\%，各类别 F1 分别为：DoS=88\%、Normal=82\%、Probe=75\%、R2L=57\%、U2R=56\%。R2L 和 U2R 因训练样本极度稀缺，F1 显著低于 DoS 和 Normal，但相比二分类中验证-测试的巨大落差，5 分类的整体表现更为稳定。

图12 展示了 MLP 在二分类测试集上的混淆矩阵。可以看到，模型对 Normal 类的识别精度很高（实际 9711 条 Normal 中被正确预测为 Normal 的占绝大多数），但攻击侧的假阴性比例明显——大量攻击样本被误判为 Normal，这是当前 0.5777 召回率的直接来源。值得注意的是，模型的高 precision（0.9560）意味着被预测为异常的结果中误报极少，但代价是漏报大量真实攻击，反映出在入侵检测这种"宁误勿漏"需求下的决策阈值尚未最优化。

\begin{figure}[htbp]
  \centering
  \includegraphics[width=\columnwidth]{09_confusion_matrix_mlp.png}
  \caption{MLP 二分类混淆矩阵}
  \label{fig:09_confusion_matrix_mlp}
\end{figure}

图13 进一步给出 MLP 多分类（5 大类）的混淆矩阵。在 5 类设定下，DoS 与 Normal 类因样本充足且特征分布清晰，识别精度较高（DoS 6069/7460、Normal 9299/9711 被正确分类）。R2L 类有较多被误判为 Normal（1626/2885 漏报），U2R 类因训练样本极少（仅 52 条）大量漏报，仅 36/67 被正确分类。这一结构清晰呈现了"样本量决定识别能力"的级联效应。

\begin{figure}[htbp]
  \centering
  \includegraphics[width=\columnwidth]{09_confusion_matrix_mlp_multiclass.png}
  \caption{MLP 五分类混淆矩阵}
  \label{fig:09_confusion_matrix_mlp_multiclass}
\end{figure}

\subsection{CNN一维卷积网络}

一维卷积网络的架构为：Conv1d(1→16, kernel\_size=3) → ReLU → MaxPool1d(2) → Flatten → Linear(2)。输入将 20 维特征视为长度 20 的序列，卷积核沿特征维度滑动以提取局部模式。

二分类测试集结果：acc=0.7635，f1=0.7452，AUC=0.8820。整体性能略优于 LSTM，但与 MLP 处于相近水平。1D CNN 的优势在具有局部相关性的信号数据（如音频、时序传感器读数）中较为明显，而 NSL-KDD 特征之间并无明显的空间局部结构，每个特征独立编码网络连接的某一统计量，因此卷积核难以捕捉到有意义的局部模式。这一结果印证了"模型架构必须匹配数据内在结构"的原则。

\subsection{LSTM长短期记忆网络}

LSTM 的架构为：LSTM(input\_size=20, hidden\_size=32) → Linear(2)，仅含一层单向 LSTM，输出最后时间步的隐藏状态送入全连接层做分类。

二分类测试集结果：acc=0.7427，f1=0.7161，AUC=0.8832。性能与 CNN 相近，但弱于 MLP。LSTM 擅长捕捉时序依赖，而 NSL-KDD 的每条样本代表一次独立的网络连接，样本之间无时序关系，因此 LSTM 的循环结构并未带来收益，反而因参数量更大、训练更慢而处于劣势。这一结果与 CNN 的讨论相互印证：对于样本独立、结构扁平的表格数据，简单的前馈架构（MLP）反而是最合适的选择。

\subsection{神经网络方法总结}

综合来看，MLP 在验证集上显著优于传统机器学习方法，AUC 接近 1.0。但在统一测试集上，神经网络与传统方法的差距大幅缩小，MLP 测试集 F1 仅为 0.7202，弱于决策树的 0.76。CNN 在测试集上的 F1（0.745）和 Accuracy（0.764）反而略高于 MLP，但优势有限且 AUC 明显低于 MLP，整体仍处于相近水平。这一对比揭示了一个常被忽视的事实：验证集上的高指标可能源于训练集与验证集分布的相似性，而测试集中的 unseen 攻击才是真实部署场景的合理近似。三种深度模型在测试集上的性能差距远小于验证集上暗示的优势，进一步佐证了"模型复杂度不等于性能提升"的原则。

在第 5 章中，将对所有 10 个模型在统一测试集下的指标进行系统对比，量化分析传统机器学习与深度学习方法在入侵检测任务上的真实差距，并结合训练耗时、可解释性、部署成本等维度给出综合的选型建议。
